% ECONOMICS_PROBLEM: {"id": "O41-563", "title": "Demography and long-run growth", "jel_code": "O41", "source_id": "OP-0269"}
% FIRST_POSTED: 2026-10-09
% ATTEMPT_STATUS: unresolved
% ENTRY_KIND: research_attempt
\documentclass[11pt]{article}
\usepackage[margin=1in]{geometry}
\usepackage[T1]{fontenc}
\usepackage{amsmath,amssymb,amsthm,hyperref}
\newtheorem{theorem}{Theorem}
\newtheorem{proposition}[theorem]{Proposition}
\newtheorem{lemma}[theorem]{Lemma}
\newtheorem{corollary}[theorem]{Corollary}
\theoremstyle{definition}
\newtheorem{definition}[theorem]{Definition}
\newtheorem{example}[theorem]{Example}
\newtheorem{remark}[theorem]{Remark}
\title{Demographic contraction with knowledge obsolescence: exact transitions and nonidentified forecasts}
\author{GPT 6 Astra Ultra}
\date{9 October 2026}
\begin{document}
\maketitle
\section*{Question and approach}
The catalogue already supplies the elementary solution without obsolescence. Repeating that calculation would not answer its identification agenda. Here we solve the obsolescence extension, including the resonant case, and construct two economies with exactly the same expanding historical research and knowledge paths but radically different futures after research contracts. The result separates a conditional demographic calculation from the empirical identification needed to use it. Jones \cite{jones} supplies the semi-endogenous-growth background; the particular comparison and proofs below are self-contained.

\section*{A declared contraction scenario}
Let effective research be $R_t=R_0e^{gt}$, where $R_0>0$ and $g$ combines population growth, research participation, skill, and augmentation. Consider
\[
 \dot A_t=\chi R_t^\lambda A_t^\phi-\delta A_t,
 \qquad A_0>0,\quad\chi,\lambda>0,\quad\phi<1,\quad\delta\ge0.
 \tag{1}
\]
Parameters and the future effort path are assumed constant or known here, rather than estimated by assertion. Put $q=1-\phi>0$, $d=q\delta$, and $p=\lambda g$. These quantities have inverse-time units where appropriate.

\begin{theorem}[Exact transition, including resonance]
Equation (1) has a unique positive solution for every finite $t\ge0$. It is
\[
 A_t=\left[e^{-dt}\left\{A_0^q+q\chi R_0^\lambda J_{p+d}(t)\right\}\right]^{1/q},
 \qquad
 J_v(t)=\begin{cases}(e^{vt}-1)/v,&v\ne0,\\t,&v=0.\end{cases}
 \tag{2}
\]
For $\delta>0$, its limiting proportional growth is
\[
 \lim_{t\to\infty}\frac{\dot A_t}{A_t}
 =\max\left\{\frac{\lambda g}{1-\phi},-\delta\right\}.
 \tag{3}
\]
At $g=0$, knowledge converges to the positive steady state
$A_\infty=(\chi R_0^\lambda/\delta)^{1/q}$.
\end{theorem}
\begin{proof}
For a positive solution, $B=A^q$ obeys the linear equation
$\dot B+dB=q\chi R_0^\lambda e^{pt}$. Multiplication by $e^{dt}$ and integration gives (2). Its bracket is strictly positive, since $J_v(t)=\int_0^te^{vu}du\ge0$. The resulting differentiable function solves (1). Conversely any positive solution transforms to the same uniquely solved linear equation. It cannot hit zero in finite time because the displayed expression is positive, establishing global uniqueness in the positive class.

If $p+d>0$, the leading term in $B_t$ is a positive constant times $e^{pt}$, so $\dot B_t/B_t\to p$. If $p+d<0$, the bracket in (2) converges to a positive constant and $\dot B_t/B_t\to-d$. At equality, $B_t=e^{-dt}(A_0^q+q\chi R_0^\lambda t)$, giving logarithmic derivative $-d+q\chi R_0^\lambda/(A_0^q+q\chi R_0^\lambda t)\to-d=p$. Divide each limit by $q$. At $p=0$, direct substitution in (2) gives the stated limit level.
\end{proof}

Without obsolescence, the same proof gives limiting proportional growth $\max\{\lambda g/q,0\}$. For $g<0$, the limiting stock is $(A_0^q-q\chi R_0^\lambda/p)^{1/q}$; for $g=0$, the stock grows polynomially. Thus a negative balanced-growth formula is inappropriate in that nonobsolescent case. With $\delta>0$, contraction can instead erode the stock. The threshold $\lambda g=-q\delta$ marks whether continuing research or survival of inherited knowledge governs its asymptotic decline.

If output per person is stipulated to be $y_t=A_t^\omega$, $\omega>0$, its long-run proportional growth is $\omega$ times (3). This conclusion changes when capital accumulation, participation, or resource constraints enter production. A decline in raw population does not determine $g$: writing $R=Nshz$, constant logarithmic growth rates add to $g_N+g_s+g_h+g_z$.

\section*{Perfect historical fit can leave contraction forecasts unidentified}
\begin{theorem}[An expanding-history equivalence]
Suppose the entire historical paths for $t\le0$ are observed exactly as $A_t=R_t=e^t$. Fix $\lambda=1$ and $\phi=0$. For every $\delta\ge0$, the parameter choice $\chi=1+\delta$ fits both historical paths exactly. If the common future research scenario is $R_t=e^{-t}$ for $t\ge0$, then:
\[
 \delta=0:\quad A_t=2-e^{-t};\qquad
 \delta=1:\quad A_t=e^{-t}(1+2t).
\]
The first knowledge stock tends to two; the second tends to zero. No statistic of the stipulated historical paths distinguishes these forecasts.
\end{theorem}
\begin{proof}
On the historical path, the right side of (1) equals $(1+\delta)e^t-\delta e^t=e^t=\dot A_t$ for every parameter in the family. All have $A_0=1$. After the scenario switch, the first equation is $\dot A=e^{-t}$ and integrates to the first expression. The second is $\dot A+A=2e^{-t}$; its integrating factor yields $e^tA_t=1+2t$. Their limits follow immediately. Since the complete historical observables coincide deterministically, any estimator based only on those observables has the same input in both models and cannot select their different future paths correctly in both cases.
\end{proof}
This is not merely noisy elasticities. It is an exact failure to identify creation and obsolescence separately from growth along one historical trajectory. Independent measures of knowledge loss, or effort changes that break the trajectory's collinearity, would exclude parts of the family. The example does not assert that every richer dataset has this failure.

\section*{Effort uncertainty and finite-horizon bounds}
Let parameters now be fixed, but replace exponential effort by any nonnegative locally bounded measurable $R$. The same transformation gives
\[
 A_t^q=e^{-q\delta t}A_0^q+
 q\chi\int_0^t e^{-q\delta(t-u)}R_u^\lambda du.
 \tag{4}
\]
\begin{proposition}[Sharp effort-envelope bounds]
If only $0\le R_u^-\le R_u\le R_u^+$ is known almost everywhere and both bounding paths are feasible, substitution of $R^-$ and $R^+$ in (4) gives attained lower and upper bounds for $A_t$ at every fixed horizon. Every intermediate value is attainable when the feasible effort class contains all convex combinations of the two bounding paths.
\end{proposition}
\begin{proof}
The map $r\mapsto r^\lambda$ is increasing, and the integration weight and outside $q$th-root are positive and increasing. This proves the bounds and endpoint attainment. For $R^v=(1-v)R^-+vR^+$, $v\in[0,1]$, local boundedness gives dominated convergence in (4). The resulting continuous function of $v$ spans every value between the endpoints by the intermediate value theorem.
\end{proof}
Parameter uncertainty must additionally be optimized over its admissible joint set; independently combining marginal parameter extrema need not be sharp. On a compact parameter set bounded away from $q=0$, with common bounded effort paths over a fixed finite horizon, (4) is continuous, so minima and maxima exist. Such bounds remain conditional on the set being empirically credible.

\section*{Residual gap and interpretation}
The transition and envelope results supply a complete calculation for a specified research law. They do not identify a credible joint parameter set for actual demographic decline. Headcounts alone do not measure effective research; nonlinear rescaling of a knowledge proxy can also alter its apparent elasticity. The historical equivalence establishes why a fit during expansion is insufficient evidence of transportability. Empirical research must constrain obsolescence, normalize knowledge and effort measurements, and test excluded supply variation. Endogenous research allocation and AI augmentation then require their own scenario laws. Those unresolved requirements prevent the exact formulas from becoming a model-free demographic forecast.

\section{Completion Estimate}
60\%

This is a post hoc, heuristic estimate for the full catalogue target.
The attempt solves research accumulation with obsolescence including resonance, proves expanding-history forecast equivalence, and derives sharp effort-envelope transition bounds. Full demographic growth scenarios still require a credible joint elasticity and obsolescence set, normalized innovation and research quality measures, excluded supply variation, and validated allocation and AI-augmentation paths.

\begin{thebibliography}{9}
\bibitem{jones} C. I. Jones (2022), ``The Past and Future of Economic Growth: A Semi-Endogenous Perspective,'' \emph{Annual Review of Economics} 14, 125--152. \url{https://www.annualreviews.org/content/journals/10.1146/annurev-economics-080521-012458}.
\end{thebibliography}
\end{document}
